\section{LightRAG 原理}

\subsection{概述}

LightRAG 是一篇 EMNLP 2025 的学术论文，也是 Paper2Slides 作者的另一个项目。它是一个基于\textbf{知识图谱}的 RAG 系统，相比传统向量 RAG 具有更强的语义理解能力。

\begin{importantbox}{LightRAG 的核心特性}
\begin{itemize}
    \item 基于 \textbf{Knowledge Graph} 的索引（而非纯向量索引）
    \item \textbf{双层检索}：Low-Level（局部，特定实体）+ High-Level（全局，主题和概念）
    \item \textbf{增量更新}：新文档无需重建整个 Graph，只需执行提取步骤并合并
\end{itemize}
\end{importantbox}

\subsection{Indexing 流程：构建知识图谱}

对每个文档执行以下步骤：

\begin{enumerate}
    \item \textbf{Chunking}：将文档分块
    \item \textbf{Entity \& Relationship Extraction}：基于 LLM 从每个 Chunk 提取实体和关系
    \item \textbf{Profiling}：基于 LLM 对实体和关系生成文本描述
    \item \textbf{Deduplication}：对整个知识图谱做去重合并
    \item \textbf{Embedding}：对所有实体和关系的 Key 构建 Embedding Vector
\end{enumerate}

知识图谱中的数据结构：
\begin{itemize}
    \item \textbf{Entity}：名字、类型、描述（Profiling 生成）、原始 Chunk ID
    \item \textbf{Edge}：Source、Target、关键词、描述、原始 Chunk ID
\end{itemize}

\begin{knowledgebox}{LightRAG vs. GraphRAG}
LightRAG 相比 GraphRAG 的一个重要优势是\textbf{增量更新}——新文档只需执行相同的提取步骤，合并实体和边即可。而 GraphRAG 需要重新进行社区发现，代价高昂。
\end{knowledgebox}

\subsection{Retrieval 流程：双层检索}

\begin{enumerate}
    \item 基于 Query，LLM 生成 \textbf{Low-Level Keywords}（特定实体）和 \textbf{High-Level Keywords}（主题/概念）
    \item 基于 Embedding 相似度在知识图谱中匹配相关节点
    \item 通过\textbf{一跳邻居}拓展 Subgraph
    \item 获取三种类型的 Context：实体、关系、原始 Chunk
    \item Context + Query $\rightarrow$ LLM 生成最终回答
\end{enumerate}

\subsection{本章小结}

LightRAG 通过知识图谱索引 + 双层检索实现了比纯向量 RAG 更强的语义理解能力，且支持增量更新。

